\begin{result}{C.1.4}{An explicit polarized pair outside the abelian image}{Proved}
Put $a=\sqrt[3]{2}>0$, $V=\Q e_1\oplus\Q e_2$, and
\[
 V^{3,0}=\C(e_1+ia e_2),\qquad
 V^{0,3}=\C(e_1-ia e_2),\qquad f=\operatorname{diag}(1,-1).
\]
This is a polarizable rational Hodge structure with a rational anti-Hodge involution. The pair $(V,f)$ is outside the essential image of the abelian-motive category in C.1.1.
\end{result}
\begin{proof}
For any positive real $a$, the displayed two lines are complementary and conjugate, and $f$ exchanges them. In the convention $h(z)|_{V^{p,q}}=z^{-p}\bar z^{-q}$ the Weil operator and the rational alternating coefficient of a polarization are
\[
 C_a=\begin{pmatrix}0&a^{-1}\\-a&0\end{pmatrix},\qquad
 E=\begin{pmatrix}0&-1\\1&0\end{pmatrix},\qquad
 q(x,y)=(2\pi i)^{-3}x^tEy.
\]
The pairing has target $\Q(-3)$, is zero on each of the two equal-type lines, and is therefore a Hodge morphism. Moreover
\[
 EC_a=\operatorname{diag}(a,a^{-1})>0,\qquad
 f^tEf=-E,\qquad fC_af=C_a^{-1}.
\]
Thus $q$ is a polarization, also compatible with the involution on its Tate target. These identities are exact, and apply to both examples in this subsection.

For $a=\sqrt[3]{2}$, put $\tau=ia$. An endomorphism $\begin{psmallmatrix}r&s\\t&u\end{psmallmatrix}$ with rational entries preserves $V^{3,0}$ precisely when
\[
 t+(u-r)\tau-s\tau^2=0.
\]
Here $\tau^2=-\sqrt[3]{4}$ has degree three over $\Q$, and $\tau$ is not real; hence $\tau$ has degree six. The displayed relation forces $s=t=0$ and $u=r$, proving $\End_{\mathrm{HS}}(V)=\Q$. A commutative Mumford--Tate group would have its rational Lie algebra inside this endomorphism algebra. Its dimension is at least two, since its real points contain the nontrivial weight scaling and Hodge circle. Their tangent directions are linearly independent. This contradicts the one-dimensional endomorphism algebra. Thus $V$ is not of CM type.

Totaro's Corollary 4.2, proved from Theorem 4.1 in \cite[p.~18]{TotaroHodge2015}, excludes this $V$ from the tensor category generated by $H^1$ of curves. The realization functor from C.1.1, followed by forgetting $f$, has image in exactly that tensor category: an Artin summand becomes a sum of type $(0,0)$ lines. This explicit composite functor proves the claimed exclusion for the pair.
\end{proof}

\begin{result}{C.1.5}{A fully specified polarized pair outside the geometric image}{Proved}
Fix the following effective enumeration $(T_n)_{n\geq1}$ of deterministic Turing machines. The tape alphabet is $\{0,1,\#,\square\}$ in that order. For each $r\geq1$ use states $0,\ldots,r$, with initial state $0$ and halt state $r$; list the tables for the $4r$ nonhalting state/symbol pairs in lexicographic order, each instruction being an output symbol, a next state, and a move from the ordered list left, stay, right. Order first by $r$. Input $n$ is the binary expansion of $n$ between two $\#$ symbols on an otherwise blank two-sided tape, with the head at the first binary digit. This specifies a single enumeration of all machines of this model. Let
\[
 H=\{n:T_n\text{ halts on input }n\},\qquad
 a_H=1+\sum_{n=1}^{\infty}\frac{\mathbf1_H(n)}{4^n}.
\]
Replace $a$ in C.1.4 by $a_H$. The resulting polarized rational Hodge pair is not the realization of a pure motive defined by algebraic correspondences over $\C$, and consequently is not the realization of such a motive over $\R$. This is an exclusion from the geometric essential image; a description of every pair in that image remains an obligation.
\end{result}
\begin{proof}
We first check the specified real number. It lies in $[1,4/3]$. The set $H$ is not decidable: a machine which, on input $n$, halts exactly when a proposed decider says $T_n(n)$ does not halt gives a contradiction at its own index. If $a_H$ were computable, its base-four digits in $\{0,1\}$ would decide $H$. More explicitly, after the first $n-1$ digits have been recovered, multiply the remaining tail by $4^n$. It lies in $[0,1/3]$ when the $n$th digit is zero, and in $[1,4/3]$ when that digit is one. A computable approximation of error less than $1/6$ separates the intervals. Induction would decide every digit, a contradiction. Thus $a_H$ is noncomputable. The polarization and involution identities of C.1.4 apply because $a_H>0$; no numerical approximation is used.

We spell out why a geometric realization would force this number to be computable. A finite presentation of a pure algebraic-correspondence motive over $\C$ uses finitely many smooth projective varieties, twists, rational projectors and cycles. All of this data descends to a finitely generated extension of $\Qb$. Spreading the defining equations and the finitely many cycles gives, after restricting a nonempty open set, smooth projective families over a smooth integral finite-type $\Qb$-scheme $S$. Cycle actions are horizontal endomorphisms of their rational cohomology local systems. Idempotence, satisfied at the given complex point, holds on its connected component, so their images form a rational variation of Hodge structures. Its Hodge numbers are constant: the image of an idempotent on each Hodge bundle has locally constant rank.

For the asserted rank-two weight-three image one has $F^3=F^2=F^1$. Griffiths transversality gives $\nabla F^3\subset\Omega_S^1\otimes F^2=\Omega_S^1\otimes F^3$. Hence $F^3$ is horizontal; its conjugate is horizontal too. Parallel transport is therefore an isomorphism of rational Hodge structures along any path in the connected component. That component contains a $\Qb$-point: such points are analytically dense in a smooth variety over $\Qb$, as follows in an etale coordinate chart from the analytic inverse-function theorem and density of $\Qb$ in $\C$. A path to that point produces a rational Hodge isomorphism with a motive over $\Qb$. This is the explicit spreading argument behind the gap-in-Hodge-numbers phenomenon noted in \cite[Introduction, pp.~1--2]{TotaroHodge2015}.

For a motive over $\Qb$, choose algebraic de Rham bases, and a rational Betti basis. Comparison entries are periods over algebraic numbers, with integer powers of $2\pi i$ for the twists. Such periods have computable real and imaginary parts by \cite[Theorem 18, p.~11]{YoshinagaPeriods2008}. Algebraic numbers are computable, and the computable complex numbers form a field: reciprocal approximation follows by first finding a positive rational lower bound for the modulus of the nonzero denominator. Consequently the coordinates of the one-dimensional $F^3$ line, and its ratio in every rational basis, are computable. Rational change of basis acts on the ratio by an explicitly invertible fractional-linear transformation with rational coefficients and preserves this property. Applied to the line $\C(1,ia_H)$, this makes $a_H$ computable, contradicting the first paragraph.

This proof uses algebraic projectors in the geometric category. It does not replace them by arbitrary Hodge projectors over an arbitrary complex coefficient field. The forgetful map from real motives to their complex base changes transfers the exclusion in the required direction.
\end{proof}

